% Preámbulo común para las soluciones de los prácticos de Programación 3
\documentclass[11pt,a4paper]{article}

\usepackage[utf8]{inputenc}
\usepackage[T1]{fontenc}
\usepackage{lmodern}
\usepackage[spanish,es-noshorthands]{babel}
\usepackage[margin=2.3cm]{geometry}
\usepackage{amsmath,amssymb,amsthm}
\usepackage{enumitem}
\usepackage{xcolor}
\usepackage{tikz}
\usetikzlibrary{arrows.meta,positioning,shapes,calc}
\usepackage[noend]{algpseudocode}
\usepackage{algorithm}
\usepackage{booktabs}
\usepackage{tcolorbox}
\usepackage[hidelinks]{hyperref}

\floatname{algorithm}{Algoritmo}
\renewcommand{\algorithmicrequire}{\textbf{Entrada:}}
\renewcommand{\algorithmicensure}{\textbf{Salida:}}
\algrenewcommand\algorithmicwhile{\textbf{mientras}}
\algrenewcommand\algorithmicdo{\textbf{hacer}}
\algrenewcommand\algorithmicfor{\textbf{para}}
\algrenewcommand\algorithmicforall{\textbf{para todo}}
\algrenewcommand\algorithmicif{\textbf{si}}
\algrenewcommand\algorithmicthen{\textbf{entonces}}
\algrenewcommand\algorithmicelse{\textbf{sino}}
\algrenewcommand\algorithmicreturn{\textbf{retornar}}
\algrenewcommand\algorithmicfunction{\textbf{función}}
\algrenewcommand\algorithmicprocedure{\textbf{procedimiento}}

\theoremstyle{plain}
\newtheorem{lema}{Lema}[section]
\newtheorem{teorema}[lema]{Teorema}
\newtheorem{proposicion}[lema]{Proposición}
\theoremstyle{definition}
\newtheorem{definicion}[lema]{Definición}
\theoremstyle{remark}
\newtheorem*{observacion}{Observación}

% Caja para el enunciado de cada ejercicio
\newtcolorbox{enunciado}{colback=gray!6,colframe=gray!50,boxrule=0.4pt,arc=2pt,left=6pt,right=6pt,top=4pt,bottom=4pt,fontupper=\small}
% Nota que vincula el ejercicio con parciales/exámenes
\newcommand{\examref}[1]{\par\noindent{\small\color{blue!50!black}\textbf{Relación con parciales/exámenes:} #1}\par\medskip}

\newcommand{\NP}{\mathsf{NP}}
\newcommand{\Pclass}{\mathsf{P}}
\newcommand{\Oh}{\mathcal{O}}

\setlength{\parskip}{0.4em}
\setlength{\parindent}{0pt}


\tikzset{
  nodo/.style={circle,draw,minimum size=7mm,inner sep=0pt},
  ar/.style={-{Stealth[length=2.5mm]},thick},
  fw/.style={-{Stealth[length=2.5mm]},thick,black},
  bw/.style={-{Stealth[length=2.5mm]},thick,red!70!black,dashed},
  cam/.style={-{Stealth[length=2.5mm]},line width=1.6pt,blue!70!black},
  lab/.style={font=\small,fill=white,inner sep=1pt}
}
% Coloca los nodos de la red del Ejercicio 1
\newcommand{\nodosred}{%
  \node[nodo] (s) at (0,0) {$s$};
  \node[nodo] (u) at (2.6,1.5) {$u$};
  \node[nodo] (w) at (2.6,0) {$w$};
  \node[nodo] (y) at (2.6,-1.5) {$y$};
  \node[nodo] (v) at (5.2,1.5) {$v$};
  \node[nodo] (x) at (5.2,0) {$x$};
  \node[nodo] (z) at (5.2,-1.5) {$z$};
  \node[nodo] (t) at (7.8,0) {$t$};
}

\title{Práctico 10 --- Flujo en redes}
\author{Programación 3 --- FING, UdelaR}
\date{}

\begin{document}
\maketitle

\paragraph{Resultados del libro que usamos.} Para una red de flujo $G=(V,E)$ con fuente $s$, sumidero $t$ y capacidades $c_e$:
\begin{itemize}
\item Para todo flujo $f$ y todo corte $s$-$t$ $(A,B)$: $v(f)=f^{\mathrm{out}}(A)-f^{\mathrm{in}}(A)\le c(A,B)$ (7.6 y 7.8).
\item \textbf{Flujo máximo -- corte mínimo (7.13):} el valor máximo de un flujo es igual a la capacidad mínima de un corte. Si $f$ no admite camino $s$-$t$ en el grafo residual $G_f$, entonces $f$ es máximo, y el conjunto $A$ de vértices alcanzables desde $s$ en $G_f$ define un corte $(A,V\setminus A)$ de capacidad mínima con $c(A,V\setminus A)=v(f)$ (7.9, 7.11).
\item \textbf{Integralidad (7.14):} si las capacidades son enteras, Ford-Fulkerson devuelve un flujo máximo con todos sus valores $f(e)$ enteros.
\item \textbf{Complejidad (7.5):} con capacidades enteras, Ford-Fulkerson se ejecuta en $\Oh(mC)$, donde $C=\sum_{e \text{ sale de } s}c_e$ (cota de la cantidad de aumentos) y $m$ acota el costo de construir $G_f$ y buscar en él un camino de aumento.
\end{itemize}

%======================================================================
\section{Ejercicio 1}
\begin{enunciado}
Sea la red de flujo $G$ con capacidad constante de valor uno y flujo-$(s,t)$ $f$ según aristas indicado en la figura.
\begin{enumerate}[label=(\alph*)]
\item Encuentre un corte $s$-$t$ de capacidad 3. Determine por inspección si ese corte es de capacidad mínima.
\item Calcule el valor del flujo. Analice si el flujo es de valor máximo.
\item Si el flujo no es de valor máximo aplique el algoritmo de Ford-Fulkerson para determinar uno. [Mostrar los pasos del algoritmo mediante secuencias de grafos residuales.]
\end{enumerate}
\end{enunciado}

\examref{Ejercicio de aplicación de Ford-Fulkerson y cortes; similar al Segundo parcial 2025, Ej.~1(c)--(d) (grafo residual, flujo máximo, corte mínimo a partir de $G_r$) y al Examen diciembre 2018, Ej.~4 (obtener un corte mínimo a partir de $G_f$).}

La red, con el flujo $f(e)$ de cada arista (todas las capacidades valen $1$):

\begin{center}
\begin{tikzpicture}
\nodosred
\draw[ar] (s) -- node[lab]{$1$} (u);
\draw[ar] (s) -- node[lab]{$0$} (w);
\draw[ar] (s) -- node[lab]{$0$} (y);
\draw[ar] (u) -- node[lab]{$0$} (v);
\draw[ar] (u) -- node[lab,pos=0.3]{$1$} (x);
\draw[ar] (w) -- node[lab,pos=0.3]{$0$} (x);
\draw[ar] (w) -- node[lab,pos=0.3]{$0$} (z);
\draw[ar] (y) -- node[lab]{$0$} (z);
\draw[ar] (v) -- node[lab]{$0$} (t);
\draw[ar] (x) -- node[lab]{$1$} (t);
\draw[ar] (z) -- node[lab]{$0$} (t);
\end{tikzpicture}
\end{center}

Aristas: $(s,u)$, $(s,w)$, $(s,y)$, $(u,v)$, $(u,x)$, $(w,x)$, $(w,z)$, $(y,z)$, $(v,t)$, $(x,t)$, $(z,t)$, con $f=1$ en $(s,u),(u,x),(x,t)$ y $f=0$ en el resto.

\subsection*{(a) Un corte de capacidad 3}

Sea $A=\{s,u,v,w,x,y,z\}$, $B=\{t\}$. Las aristas que salen de $A$ hacia $B$ son $(v,t),(x,t),(z,t)$, así que
\[ c(A,B)=1+1+1=3. \]
(Otro corte de capacidad 3 es $(\{s\},V\setminus\{s\})$, con aristas salientes $(s,u),(s,w),(s,y)$.)

\textbf{Es de capacidad mínima.} Por inspección, la red tiene tres caminos $s$-$t$ que no comparten aristas:
\[ s\to u\to v\to t,\qquad s\to w\to x\to t,\qquad s\to y\to z\to t. \]
Enviando una unidad por cada uno se obtiene un flujo válido (cada arista lleva a lo sumo $1$ y se conserva el flujo en cada nodo intermedio) de valor $3$. Como el valor de cualquier flujo es menor o igual que la capacidad de cualquier corte, todo corte tiene capacidad $\ge 3$. Luego el corte $(A,B)$ de capacidad $3$ es de capacidad mínima, y el flujo máximo vale $3$.

\subsection*{(b) Valor del flujo $f$}

$v(f)=f^{\mathrm{out}}(s)=f(s,u)+f(s,w)+f(s,y)=1+0+0=1$. (Se verifica que $f$ es un flujo: cumple las capacidades y la conservación; en $u$ entra $1$ y sale $1$ por $(u,x)$, en $x$ entra $1$ y sale $1$ por $(x,t)$, y en el resto de los nodos intermedios entra y sale $0$.)

$f$ \textbf{no es máximo}: en (a) mostramos un flujo de valor $3>1$. Equivalentemente, $v(f)=1<3=$ capacidad mínima de un corte.

\subsection*{(c) Ford-Fulkerson a partir de $f$}

Como todas las capacidades son $1$ y los flujos son $0$ o $1$, en el grafo residual cada arista $e=(a,b)$ de $G$ aparece exactamente de una forma, con capacidad residual $1$: como arista \emph{hacia adelante} $(a,b)$ si $f(e)=0$, o como arista \emph{hacia atrás} $(b,a)$ si $f(e)=1$. En las figuras, las aristas hacia adelante son negras y las hacia atrás rojas punteadas (todas de capacidad residual $1$); en azul grueso se marca el camino de aumento elegido.

\paragraph{Iteración 1.} Grafo residual $G_f$ del flujo inicial ($v(f)=1$). Elegimos el camino de aumento $P_1: s\to y\to z\to t$, de cuello de botella $1$.

\begin{center}
\begin{tikzpicture}
\nodosred
\draw[bw] (u) -- (s);
\draw[fw] (s) -- (w);
\draw[cam] (s) -- (y);
\draw[fw] (u) -- (v);
\draw[bw] (x) -- (u);
\draw[fw] (w) -- (x);
\draw[fw] (w) -- (z);
\draw[cam] (y) -- (z);
\draw[fw] (v) -- (t);
\draw[bw] (t) -- (x);
\draw[cam] (z) -- (t);
\node at (3.9,-2.4) {$G_f$, \ $v(f)=1$};
\end{tikzpicture}
\end{center}

Aumentando por $P_1$ obtenemos $f_2$ con $f_2(s,y)=f_2(y,z)=f_2(z,t)=1$ y $v(f_2)=2$.

\paragraph{Iteración 2.} Grafo residual $G_{f_2}$. Ya no hay camino ``solo hacia adelante'': la única forma de llegar a $t$ es por $(v,t)$, y para llegar a $v$ hay que pasar por $u$, a lo que solo se llega usando la arista hacia atrás $(x,u)$. Elegimos $P_2: s\to w\to x\to u\to v\to t$, de cuello de botella $1$.

\begin{center}
\begin{tikzpicture}
\nodosred
\draw[bw] (u) -- (s);
\draw[cam] (s) -- (w);
\draw[bw] (y) -- (s);
\draw[cam] (u) -- (v);
\draw[cam,red!70!black] (x) -- (u);
\draw[cam] (w) -- (x);
\draw[fw] (w) -- (z);
\draw[bw] (z) -- (y);
\draw[cam] (v) -- (t);
\draw[bw] (t) -- (x);
\draw[bw] (t) -- (z);
\node at (3.9,-2.4) {$G_{f_2}$, \ $v(f_2)=2$};
\end{tikzpicture}
\end{center}

Al aumentar por $P_2$: $f(s,w)$, $f(w,x)$, $f(u,v)$, $f(v,t)$ pasan de $0$ a $1$, y, como $(x,u)$ es arista hacia atrás, $f(u,x)$ pasa de $1$ a $0$ (se ``deshace'' el flujo que iba por $u\to x$ y se redirige por $u\to v$). Obtenemos $f_3$ con $v(f_3)=3$:

\begin{center}
\begin{tikzpicture}
\nodosred
\draw[ar] (s) -- node[lab]{$1$} (u);
\draw[ar] (s) -- node[lab]{$1$} (w);
\draw[ar] (s) -- node[lab]{$1$} (y);
\draw[ar] (u) -- node[lab]{$1$} (v);
\draw[ar] (u) -- node[lab,pos=0.3]{$0$} (x);
\draw[ar] (w) -- node[lab,pos=0.3]{$1$} (x);
\draw[ar] (w) -- node[lab,pos=0.3]{$0$} (z);
\draw[ar] (y) -- node[lab]{$1$} (z);
\draw[ar] (v) -- node[lab]{$1$} (t);
\draw[ar] (x) -- node[lab]{$1$} (t);
\draw[ar] (z) -- node[lab]{$1$} (t);
\node at (3.9,-2.4) {Flujo $f_3$ en $G$, \ $v(f_3)=3$};
\end{tikzpicture}
\end{center}

\paragraph{Iteración 3.} Grafo residual $G_{f_3}$: las tres aristas que salen de $s$ están saturadas, así que en $G_{f_3}$ el nodo $s$ no tiene aristas salientes y no hay camino de $s$ a $t$. El algoritmo termina.

\begin{center}
\begin{tikzpicture}
\nodosred
\draw[bw] (u) -- (s);
\draw[bw] (w) -- (s);
\draw[bw] (y) -- (s);
\draw[bw] (v) -- (u);
\draw[fw] (u) -- (x);
\draw[bw] (x) -- (w);
\draw[fw] (w) -- (z);
\draw[bw] (z) -- (y);
\draw[bw] (t) -- (v);
\draw[bw] (t) -- (x);
\draw[bw] (t) -- (z);
\node at (3.9,-2.4) {$G_{f_3}$, \ $v(f_3)=3$};
\end{tikzpicture}
\end{center}

\textbf{Conclusión.} $f_3$ es un flujo máximo, de valor $3$. El conjunto de vértices alcanzables desde $s$ en $G_{f_3}$ es $A^\ast=\{s\}$, y el corte $(\{s\},V\setminus\{s\})$ tiene capacidad $3=v(f_3)$, lo que certifica la optimalidad (y confirma lo visto en (a)).

%======================================================================
\section{Ejercicio 2}
\begin{enunciado}
(Kleinberg \& Tardos, Ex.~7.5). Indicar si la siguiente afirmación es verdadera o falsa. Si es verdadera, dar una breve explicación. Si es falsa, dar un contraejemplo.

\emph{Sea una red de flujo $G$ con fuente $s$, destino $t$ y capacidad positiva $c_e$ para cada arista $e\in E$. Sea $(A,B)$ un corte de mínima capacidad del flujo $(s,t)$. Si se incrementa en una unidad la capacidad de cada arista entonces el corte $(A,B)$ sigue siendo de mínima capacidad con respecto a las capacidades $\{c_e+1 : e\in E\}$ en G.}
\end{enunciado}

\examref{Examen julio 2023, Ej.~(b) del ejercicio de flujo: es exactamente esta afirmación. Usamos el mismo contraejemplo de la solución oficial.}

\textbf{La afirmación es falsa.}

\textbf{Idea.} Si $k(A,B)$ es la cantidad de aristas que salen de $A$ hacia $B$, con las nuevas capacidades
\[ c'(A,B)=\sum_{e \text{ sale de } A}(c_e+1)=c(A,B)+k(A,B). \]
Un corte mínimo con muchas aristas puede ser superado por otro corte de igual (o un poco mayor) capacidad original pero con menos aristas.

\textbf{Contraejemplo.} Sea la red con aristas $(s,a)$ de capacidad $3$, $(s,b)$ de capacidad $4$, $(a,b)$ de capacidad $1$, $(a,t)$ de capacidad $1$ y $(b,t)$ de capacidad $5$:

\begin{center}
\begin{tikzpicture}
\node[nodo] (s) at (0,0) {$s$};
\node[nodo] (a) at (2.5,1.3) {$a$};
\node[nodo] (b) at (2.5,-1.3) {$b$};
\node[nodo] (t) at (5,0) {$t$};
\draw[ar] (s) -- node[lab]{$3$} (a);
\draw[ar] (s) -- node[lab]{$4$} (b);
\draw[ar] (a) -- node[lab]{$1$} (b);
\draw[ar] (a) -- node[lab]{$1$} (t);
\draw[ar] (b) -- node[lab]{$5$} (t);
\end{tikzpicture}
\end{center}

Los cortes $s$-$t$ son los cuatro conjuntos $A\ni s$ con $t\notin A$:

\begin{center}
\begin{tabular}{lllcc}
\toprule
$A$ & aristas salientes de $A$ & $k$ & $c(A,B)$ & $c'(A,B)$ \\
\midrule
$\{s\}$ & $(s,a),(s,b)$ & 2 & $3+4=7$ & $9$ \\
$\{s,a\}$ & $(s,b),(a,b),(a,t)$ & 3 & $4+1+1=6$ & $9$ \\
$\{s,b\}$ & $(s,a),(b,t)$ & 2 & $3+5=8$ & $10$ \\
$\{s,a,b\}$ & $(a,t),(b,t)$ & 2 & $1+5=6$ & $8$ \\
\bottomrule
\end{tabular}
\end{center}

Con las capacidades originales, $(A,B)=(\{s,a\},\{b,t\})$ es un corte de capacidad mínima ($6$). Con las capacidades incrementadas su capacidad pasa a ser $9$, mientras que $(\{s,a,b\},\{t\})$ tiene capacidad $8<9$. Por lo tanto $(\{s,a\},\{b,t\})$ deja de ser de capacidad mínima, y la afirmación es falsa.

\begin{observacion}
En cambio, si todas las capacidades se \emph{multiplican} por una constante $\lambda>0$, todo corte mínimo sigue siendo mínimo, porque $c'(A,B)=\lambda\, c(A,B)$ para todo corte y el orden entre cortes se preserva. Al \emph{sumar} una constante, el término $k(A,B)$ depende del corte y el orden puede cambiar, como en el ejemplo.
\end{observacion}

%======================================================================
\section{Ejercicio 3}
\begin{enunciado}
(Kleinberg \& Tardos, Ex.~7.9). Por inundaciones, hay $n$ pacientes que deben ser trasladados a alguno de los $k$ hospitales de la región. Cada paciente debe ir a un hospital que esté a lo sumo a media hora de traslado de su ubicación (distintos pacientes tienen distintas opciones). Además, para no sobrecargar los hospitales, cada hospital debe recibir a lo sumo $\lceil n/k\rceil$ pacientes. Se quiere decidir si se puede elegir un hospital para cada paciente cumpliendo estas condiciones.
\begin{enumerate}[label=(\alph*)]
\item Proporcione un algoritmo de tiempo polinomial en $n$ y $k$ que, a partir de las ubicaciones de los pacientes, determine si el traslado es posible en dichas condiciones. Sugerencia: modele el problema mediante una red con $k$ nodos que representan hospitales, $n$ nodos que representan pacientes, un nodo fuente y un nodo sumidero.
\item Muestre la corrección de su algoritmo.
\item Analice la complejidad de su algoritmo.
\end{enumerate}
\end{enunciado}

\examref{La técnica (red bipartita fuente $\to$ oferta $\to$ demanda $\to$ sumidero, y respuesta afirmativa si y solo si el flujo máximo satura las aristas de un lado) es la de la solución oficial del Examen febrero 2023, Ej.~1 (reservas de sangre de un hospital); ver también Examen diciembre 2021, Ej.~3 (asignación de personal de enfermería).}

Suponemos que, a partir de las ubicaciones, se puede determinar en $\Oh(1)$ si el hospital $j$ está a lo sumo a media hora del paciente $i$ (por ejemplo, contando con la tabla de tiempos de traslado, o con una función que los calcula). Llamamos $R_i\subseteq\{1,\dots,k\}$ al conjunto de hospitales admisibles para el paciente $i$, y $L=\lceil n/k\rceil$.

Una \emph{asignación válida} es una función $h:\{1,\dots,n\}\to\{1,\dots,k\}$ tal que $h(i)\in R_i$ para todo $i$ y $|\{i : h(i)=j\}|\le L$ para todo $j$.

\subsection*{(a) Algoritmo}

\textbf{Construcción de la red $G=(V,E)$.}
\begin{itemize}
\item Nodos: una fuente $s$, un sumidero $t$, un nodo $p_i$ por cada paciente ($1\le i\le n$) y un nodo $q_j$ por cada hospital ($1\le j\le k$).
\item Aristas $(s,p_i)$ con capacidad $1$, para cada paciente $i$.
\item Aristas $(p_i,q_j)$ con capacidad $1$, para cada par con $j\in R_i$ (el hospital $j$ está a lo sumo a media hora del paciente $i$).
\item Aristas $(q_j,t)$ con capacidad $L=\lceil n/k\rceil$, para cada hospital $j$.
\end{itemize}

\begin{center}
\begin{tikzpicture}[xscale=1.1]
\node[nodo] (s) at (0,0) {$s$};
\node[nodo] (p1) at (2.5,1.8) {$p_1$};
\node[nodo] (p2) at (2.5,0.6) {$p_2$};
\node (pd) at (2.5,-0.5) {$\vdots$};
\node[nodo] (pn) at (2.5,-1.8) {$p_n$};
\node[nodo] (q1) at (5.5,1.2) {$q_1$};
\node (qd) at (5.5,0.1) {$\vdots$};
\node[nodo] (qk) at (5.5,-1.2) {$q_k$};
\node[nodo] (t) at (8,0) {$t$};
\draw[ar] (s) -- node[lab]{$1$} (p1);
\draw[ar] (s) -- node[lab]{$1$} (p2);
\draw[ar] (s) -- node[lab]{$1$} (pn);
\draw[ar] (p1) -- node[lab]{$1$} (q1);
\draw[ar] (p2) -- node[lab,pos=0.3]{$1$} (qk);
\draw[ar] (p2) -- node[lab,pos=0.3]{$1$} (q1);
\draw[ar] (pn) -- node[lab]{$1$} (qk);
\draw[ar] (q1) -- node[lab]{$\lceil n/k\rceil$} (t);
\draw[ar] (qk) -- node[lab]{$\lceil n/k\rceil$} (t);
\end{tikzpicture}
\end{center}

\begin{algorithm}[H]
\caption{Traslado equilibrado de pacientes}
\begin{algorithmic}[1]
\Require $n$ pacientes, $k$ hospitales, y para cada par $(i,j)$ si el hospital $j$ está a lo sumo a media hora del paciente $i$
\Ensure ``sí'' (y una asignación válida) si existe una, ``no'' en caso contrario
\State Construir la red $G$ descrita arriba
\State $f\gets$ \Call{FordFulkerson}{$G,s,t$} \Comment{flujo máximo entero}
\If{$v(f)=n$}
  \ForAll{paciente $i$}
     \State $h(i)\gets$ el único $j$ con $f(p_i,q_j)=1$
  \EndFor
  \State \Return ``sí'', $h$
\Else
  \State \Return ``no''
\EndIf
\end{algorithmic}
\end{algorithm}

\subsection*{(b) Corrección}

Probamos: \emph{existe una asignación válida si y solo si el valor del flujo máximo en $G$ es $n$}; y si es $n$, la $h$ construida en el algoritmo es válida.

Primero, todo flujo tiene valor $\le n$, porque el corte $(\{s\},V\setminus\{s\})$ tiene capacidad $\sum_{i=1}^n c(s,p_i)=n$.

\textbf{($\Rightarrow$) Si existe una asignación válida, el flujo máximo vale $n$.} Sea $h$ una asignación válida. Definimos
\[ f(s,p_i)=1,\qquad f(p_i,q_j)=\begin{cases}1 & \text{si } j=h(i)\\ 0 & \text{si no}\end{cases},\qquad f(q_j,t)=|\{i : h(i)=j\}|. \]
Las aristas $(p_i,q_{h(i)})$ existen porque $h(i)\in R_i$. Veamos que $f$ es un flujo:
\begin{itemize}
\item \emph{Capacidades:} en las aristas $(s,p_i)$ y $(p_i,q_j)$ el flujo es $0$ o $1$, y la capacidad es $1$; en $(q_j,t)$, $f(q_j,t)=|\{i:h(i)=j\}|\le L$ porque $h$ es válida.
\item \emph{Conservación en $p_i$:} entra $f(s,p_i)=1$ y sale $\sum_j f(p_i,q_j)=1$ (solo $j=h(i)$).
\item \emph{Conservación en $q_j$:} entra $\sum_i f(p_i,q_j)=|\{i:h(i)=j\}|$ y sale $f(q_j,t)$, que es lo mismo.
\end{itemize}
Su valor es $v(f)=\sum_i f(s,p_i)=n$. Como ningún flujo supera $n$, el flujo máximo vale $n$, y por lo tanto Ford-Fulkerson devuelve un flujo de valor $n$.

\textbf{($\Leftarrow$) Si el flujo máximo vale $n$, la $h$ construida es una asignación válida.} Sea $f$ el flujo devuelto por Ford-Fulkerson, con $v(f)=n$. Como las capacidades son enteras, $f$ es entero (7.14).
\begin{itemize}
\item Como $v(f)=\sum_{i=1}^n f(s,p_i)=n$ y cada $f(s,p_i)\le 1$, se tiene $f(s,p_i)=1$ para todo $i$.
\item Por conservación en $p_i$, $\sum_j f(p_i,q_j)=1$; como cada $f(p_i,q_j)$ es un entero en $\{0,1\}$, hay \emph{exactamente un} $j$ con $f(p_i,q_j)=1$. Ese es $h(i)$, que está bien definido, y $h(i)\in R_i$ porque la arista $(p_i,q_{h(i)})$ existe solo si el hospital $h(i)$ está a lo sumo a media hora del paciente $i$.
\item Para cada hospital $j$, por conservación en $q_j$: $|\{i:h(i)=j\}|=\sum_i f(p_i,q_j)=f(q_j,t)\le L=\lceil n/k\rceil$.
\end{itemize}
Luego $h$ es una asignación válida.

\textbf{Conclusión.} Si existe una asignación válida, por ($\Rightarrow$) Ford-Fulkerson devuelve un flujo de valor $n$ y el algoritmo responde ``sí'' con una asignación válida (por ($\Leftarrow$)). Si no existe, por ($\Leftarrow$) (contrarrecíproco) el flujo máximo no puede valer $n$, así que vale menos y el algoritmo responde ``no''. \qed

(Notar que si algún paciente no tiene ningún hospital a menos de media hora, $p_i$ no tiene aristas salientes, ningún flujo puede pasar por $(s,p_i)$, el máximo es $<n$ y el algoritmo responde correctamente ``no''.)

\subsection*{(c) Complejidad}

\begin{itemize}
\item \textbf{Tamaño de la red:} $|V|=n+k+2$ y $|E|=n+\sum_i|R_i|+k\le n+nk+k$. Como $k\ge1$, $m:=|E|=\Oh(nk)$ (pues $n\le nk$ y $k\le nk$).
\item \textbf{Construcción:} decidir para cada uno de los $nk$ pares $(i,j)$ si $j\in R_i$ cuesta $\Oh(1)$ cada uno, más crear nodos y aristas: $\Oh(nk)$.
\item \textbf{Ford-Fulkerson:} por (7.5), $\Oh(mC)$ con $C=\sum_{e\text{ sale de }s}c_e=n$. Cada una de las a lo sumo $n$ iteraciones construye $G_f$ y busca un camino $s$-$t$ (BFS/DFS) en $\Oh(|V|+|E|)=\Oh(nk)$. Total: $\Oh(n\cdot nk)=\Oh(n^2k)$.
\item \textbf{Construir $h$:} recorrer las aristas $(p_i,q_j)$: $\Oh(nk)$.
\end{itemize}

Tiempo total: $\Oh(n^2k)$, polinomial en $n$ y $k$.

\end{document}
